\documentclass[12pt,a4paper]{article}

\usepackage[utf8]{inputenc}
\usepackage[T1]{fontenc}
\usepackage[spanish,es-tabla]{babel}
\usepackage{geometry}
\usepackage{array}
\usepackage{longtable}
\usepackage{booktabs}
\usepackage{xcolor}
\usepackage{hyperref}
\usepackage{enumitem}

\geometry{left=2.5cm,right=2.5cm,top=2.5cm,bottom=2.5cm}
\setlength{\parindent}{0pt}
\setlength{\parskip}{6pt}
\setlist[itemize]{leftmargin=1.2cm}
\setlist[enumerate]{leftmargin=1.2cm}

\hypersetup{
  colorlinks=true,
  linkcolor=black,
  urlcolor=blue
}

\newcommand{\linea}{\rule{0.92\linewidth}{0.4pt}}
\newcommand{\pendiente}{\textcolor{red}{Pendiente}}
\newcommand{\ok}{\textcolor{green!50!black}{Definido}}
\newcolumntype{L}[1]{>{\raggedright\arraybackslash}p{#1}}

\begin{document}

\begin{center}
{\Large \textbf{Cuestionario de levantamiento de datos}}\\
{\large Capitulo I: Memoria descriptiva}\\
{\normalsize Proyecto de instalacion electrica domiciliaria - vivienda propia de 2 pisos}
\end{center}

\section*{1. Objetivo del cuestionario}

Este cuestionario sirve para reunir la informacion minima necesaria antes de redactar el Capitulo I y preparar los datos de entrada del Capitulo II. El proyecto se desarrollara sobre una vivienda unifamiliar propia de 2 pisos, no sobre la vivienda anterior de referencia.

El cuestionario tambien deja indicado que norma sustenta cada bloque de informacion. Para el arranque se usaran como base la Norma EM.010 Instalaciones Electricas Interiores y el Codigo Nacional de Electricidad - Utilizacion.

\section*{2. Normativa base revisada}

\begin{longtable}{L{3.2cm}L{4.6cm}L{5.8cm}}
\toprule
\textbf{Fuente} & \textbf{Referencia inicial} & \textbf{Uso en el proyecto} \\
\midrule
\endfirsthead
\toprule
\textbf{Fuente} & \textbf{Referencia inicial} & \textbf{Uso en el proyecto} \\
\midrule
\endhead
RNE EM.010 & Articulos 1 y 2, paginas 1--2 extraidas & Define que la instalacion electrica interior va desde la acometida hasta los puntos de utilizacion y aplica a viviendas. \\
CNE-U & Seccion 010, pagina 6 extraida & Sustenta el uso obligatorio del Codigo Nacional de Electricidad - Utilizacion. \\
CNE-U & Pagina 13 extraida & Sirve para revisar el concepto de unidad de vivienda. \\
CNE-U & Seccion 050, pagina 62 extraida & Sirve para iniciar criterios de demanda en viviendas unifamiliares. \\
CNE-U & Seccion 050, pagina 69 extraida & Sirve para revisar circuitos derivados en unidades de vivienda. \\
\bottomrule
\end{longtable}

\textbf{Evidencias guardadas:} \texttt{avance-capitulos-1-2/normativa-extraida/}

\section*{3. Datos generales del estudiante y del proyecto}

\begin{longtable}{L{6cm}L{8cm}}
\toprule
\textbf{Pregunta} & \textbf{Respuesta} \\
\midrule
\endfirsthead
\toprule
\textbf{Pregunta} & \textbf{Respuesta} \\
\midrule
\endhead
Nombre del estudiante & \linea \\
Curso & Instalaciones Electricas I \\
Docente / ingeniero & \linea \\
Fecha de avance & \linea \\
Titulo tentativo del proyecto & Instalacion electrica interior de vivienda unifamiliar de 2 pisos \\
Tipo de trabajo & Academico, no expediente ejecutivo \\
\bottomrule
\end{longtable}

\section*{4. Ubicacion y datos de la vivienda}

\begin{longtable}{L{6cm}L{8cm}}
\toprule
\textbf{Pregunta} & \textbf{Respuesta} \\
\midrule
\endfirsthead
\toprule
\textbf{Pregunta} & \textbf{Respuesta} \\
\midrule
\endhead
Direccion o referencia de la vivienda & \linea \\
Distrito & \linea \\
Provincia & \linea \\
Departamento & \linea \\
Zona & Urbana / periurbana / rural: \linea \\
Tipo de vivienda & Vivienda unifamiliar \\
Numero de pisos & 2 pisos \\
Material predominante & Ladrillo / concreto / adobe / otro: \linea \\
Area aproximada del primer piso & \linea \\
Area aproximada del segundo piso & \linea \\
Area total aproximada & \linea \\
Empresa distribuidora local & \linea \\
Tension de suministro asumida & 220 V, 60 Hz, monofasico \\
\bottomrule
\end{longtable}

\section*{5. Distribucion arquitectonica preliminar}

\subsection*{5.1 Primer piso}

\begin{longtable}{L{3.4cm}L{2.4cm}L{3.8cm}L{4cm}}
\toprule
\textbf{Ambiente} & \textbf{Focos} & \textbf{Tomacorrientes} & \textbf{Observacion} \\
\midrule
\endfirsthead
\toprule
\textbf{Ambiente} & \textbf{Focos} & \textbf{Tomacorrientes} & \textbf{Observacion} \\
\midrule
\endhead
Cuarto 1 grande & 2 & 4 & Un tomacorriente por pared \\
Cuarto 2 grande & 2 & 4 & Un tomacorriente por pared \\
Cocina & 2 & 1 & Para licuadora y pequenos artefactos \\
Bano & \pendiente & \pendiente & Confirmar si se incluye \\
Escalera / pasadizo & \pendiente & \pendiente & Confirmar si se incluye \\
\bottomrule
\end{longtable}

\subsection*{5.2 Segundo piso}

\begin{longtable}{L{3.4cm}L{2.4cm}L{3.8cm}L{4cm}}
\toprule
\textbf{Ambiente} & \textbf{Focos} & \textbf{Tomacorrientes} & \textbf{Decision pendiente} \\
\midrule
\endfirsthead
\toprule
\textbf{Ambiente} & \textbf{Focos} & \textbf{Tomacorrientes} & \textbf{Decision pendiente} \\
\midrule
\endhead
Habitacion 1 & 1 o 2 & 3 o 4 & Definir segun tamano \\
Habitacion 2 & 1 o 2 & 3 o 4 & Definir segun tamano \\
Habitacion 3 & 1 o 2 & 3 o 4 & Definir segun tamano \\
Habitacion 4 & 1 o 2 & 3 o 4 & Definir segun tamano \\
Bano & \pendiente & \pendiente & Confirmar si existe \\
Escalera / pasadizo & \pendiente & \pendiente & Confirmar puntos de paso \\
\bottomrule
\end{longtable}

\section*{6. Preguntas tecnicas para cerrar el Capitulo I}

\begin{enumerate}
  \item ¿La vivienda tendra solo los ambientes indicados o se agregara bano, lavanderia, sala, comedor o patio?
  \item ¿La cocina solo usara un tomacorriente para licuadora y pequenos artefactos, o tambien se incluira refrigeradora?
  \item ¿Habra ducha electrica, terma, bomba de agua, lavadora, horno electrico u otra carga especial?
  \item ¿El tablero general estara en el primer piso cerca del ingreso principal?
  \item ¿La instalacion sera empotrada, superficial o mixta?
  \item ¿Se dibujara un croquis propio por piso o se usara un plano arquitectonico existente?
  \item ¿El ingeniero exige medidas exactas de ambientes o basta un plano/croquis con dimensiones aproximadas?
  \item ¿Se considerara sistema de puesta a tierra desde este avance o se dejara para el desarrollo tecnico del Capitulo II?
\end{enumerate}

\section*{7. Circuitos preliminares para guiar el Capitulo II}

\begin{longtable}{L{2.2cm}L{4.5cm}L{6.8cm}}
\toprule
\textbf{Circuito} & \textbf{Uso} & \textbf{Ambientes atendidos} \\
\midrule
\endfirsthead
\toprule
\textbf{Circuito} & \textbf{Uso} & \textbf{Ambientes atendidos} \\
\midrule
\endhead
C1 & Alumbrado primer piso & Cuarto 1, cuarto 2 y cocina \\
C2 & Tomacorrientes primer piso & Cuarto 1 y cuarto 2 \\
C3 & Cocina & Tomacorriente de cocina para licuadora y artefactos menores \\
C4 & Alumbrado segundo piso & Cuatro habitaciones y posibles circulaciones \\
C5 & Tomacorrientes segundo piso & Cuatro habitaciones \\
\bottomrule
\end{longtable}

\section*{8. Datos que deben quedar listos antes de redactar}

\begin{itemize}
  \item Ubicacion real o referencia de la vivienda.
  \item Area aproximada por piso.
  \item Croquis simple del primer y segundo piso.
  \item Numero final de focos por ambiente.
  \item Numero final de tomacorrientes por ambiente.
  \item Lista de cargas especiales, si existen.
  \item Ubicacion tentativa del tablero general.
  \item Tipo de canalizacion: empotrada, superficial o mixta.
  \item Normas y paginas usadas como sustento.
\end{itemize}

\section*{9. Primeras decisiones recomendadas}

Para avanzar sin complicar el proyecto se recomienda iniciar con una vivienda modelo simple:

\begin{itemize}
  \item Primer piso: dos cuartos grandes y una cocina.
  \item Segundo piso: cuatro habitaciones.
  \item Circuitos separados para alumbrado, tomacorrientes y cocina.
  \item Sin cargas especiales en la primera version.
  \item Agregar bano, lavanderia o bomba solo si el ingeniero lo solicita o si se quiere hacer un diseno mas completo.
\end{itemize}

\section*{10. Checklist de avance}

\begin{longtable}{L{9cm}L{3.8cm}}
\toprule
\textbf{Actividad} & \textbf{Estado} \\
\midrule
\endfirsthead
\toprule
\textbf{Actividad} & \textbf{Estado} \\
\midrule
\endhead
Definir vivienda de 2 pisos como caso del proyecto & \ok \\
Registrar distribucion preliminar del primer piso & \ok \\
Registrar distribucion preliminar del segundo piso & \pendiente \\
Confirmar area aproximada por piso & \pendiente \\
Confirmar cargas especiales & \pendiente \\
Extraer evidencia normativa inicial & \ok \\
Redactar Capitulo I completo & \pendiente \\
Pasar datos al Capitulo II & \pendiente \\
\bottomrule
\end{longtable}

\end{document}
